\labsection{10}{Discover structure with clustering and PCA}{sec:lab10}

\begin{labproject}{Lab 10 unsupervised learning: discover structure without inventing meaning}
\textbf{Datasets.} \texttt{mall\_customers\_dataset.csv} and \texttt{wine\_recognition\_dataset.csv}.\\
\textbf{Scenario.} You have two unsupervised tasks: identify potentially useful customer groupings and compress correlated chemical measurements. In both cases, the algorithm returns mathematical structure; interpretation must be earned through diagnostics and domain context.

\textbf{Learning outcomes.} You should be able to standardise features for distance-based methods, assess cluster stability and choice of $k$, profile clusters in original units, fit PCA without label leakage, and interpret explained variance and loadings.

\textbf{Part A---k-means clustering.}
\begin{enumerate}
  \item Select age, annual income, and spending score; exclude identifiers. Standardise the features and explain why scale matters for Euclidean distance.
  \item Fit k-means for several candidate values of $k$ using multiple random starts. Record total within-cluster sum of squares and inspect whether an elbow is visible.
  \item Choose one $k$ as a working solution and rerun it with a large \texttt{nstart}. Report cluster sizes and centroids in the \emph{original units}.
  \item Plot income against spending score by cluster and explain what information is omitted because age is the third clustering dimension.
  \item Test stability by repeating the fit with different random seeds. Distinguish stable broad structure from assignments that change easily.
\end{enumerate}

\textbf{Part B---principal component analysis.}
\begin{enumerate}
  \item Exclude \texttt{sample\_id} and \texttt{cultivar} before PCA. Standardise all 13 chemical measurements.
  \item Report explained variance for the leading components and determine how many components are needed to capture a substantial proportion of total variance.
  \item Inspect the largest absolute loadings on PC1 and PC2 and describe which original chemical measurements drive those directions.
  \item Plot PC1 against PC2 without using cultivar labels during fitting. Introduce cultivar only afterward to assess whether known classes align with the unsupervised representation.
\end{enumerate}

\textbf{Required evidence.} Submit one cluster-selection diagnostic, one cluster profile, one PCA explained-variance summary, and one score plot. For each task, write a two-sentence interpretation that separates the algorithmic result from the substantive meaning you are assigning to it.

\textbf{Professional standard.} Do not name clusters as customer ``types'' unless the labels are supported by stable feature profiles and a real decision context. Do not describe a principal component as a latent concept unless its loading pattern and domain meaning justify that interpretation.

\tcbbreak
\lstinputlisting[style=dsr]{all_experiments/lab10_unsupervised.R}
\end{labproject}
